% new:
Having learned a set of dictionary features, we want to understand whether our learned features display reduced polysemanticity, and are therefore more interpretable. To do this in a scalable manner, we require a metric to measure how interpretable a dictionary feature is. We use the automated approach introduced in \citet{bills2023language} because it scales well to measuring interpretability on the thousands of dictionary features our autoencoders learn. In summary, the autointerpretability procedure takes samples of text where the dictionary feature activates, asks a language model to write a human-readable interpretation of the dictionary feature, and then prompts the language model to use this description to predict the dictionary feature's activation on other samples of text. The correlation between the model's predicted activations and the actual activations is that feature's interpretability score. See Appendix \ref{appendix:autointerp_protocol} and \citet{bills2023language} for further details.

% old:
% Having learned a set of dictionary features, we require a metric to measure how interpretable they are. While \citet{yun2023transformer} uses hand-counting to determine whether a dictionary feature applies to a sample of text, we instead opt to use the automated, scalable approach introduced in \citet{bills2023language}. In summary, we take samples of text where the dictionary feature activates, ask a language model to write a human-readable interpretation of the dictionary feature, and then prompt the language model to use this description to predict the dictionary feature's activation on other samples of text. The correlation between the model's predicted activations and the actual activations is that feature's interpretability score. See Appendix \ref{appendix:autointerp_protocol} and \citet{bills2023language} for further details.

% We can assess the performance of our learned dictionaries in several ways, such as by measuring their reconstruction loss or sparsity. However, these metrics do not measure progress towards our ultimate goal of interpreting language models, so we require a measure of human interpretability. 

% \citet{yun2023transformer} demonstrate the surprising level of interpretability of their learned `transformer factors' by manually looking through top-activating examples of sparse features to find patterns, and then counting the proportion of the top 200 activating cases which match the pattern. Unfortunately this method is time consuming and difficult to scale to many layers and dictionaries. Fortunately, language models themselves have recently been used to automate the interpretation of features and the evaluation of these interpretations. \citet{bills2023language} developed a protocol of for evaluating the interpretability of a neuron with a score in $[-1, 1]$ which measures the correlation over a small dataset between the true activations and the activations predicted by a language model, based on a short description of the meaning of the neuron made by a language model from seeing top-activating fragments. but apply their method to our learnt features, rather than neurons. See Appendix \ref{appendix:autointerp_protocol} and \citet{bills2023language} for further details.

We show descriptions and top-and-random scores for five dictionary features from the layer 1 residual stream in Table \ref{table:autointerpretation}. The features shown are the first five under the (arbitrary) ordering in the dictionary.

\begin{table}
    \centering
\begin{tabular}{|c|p{8cm}|c|}  
 \hline
 Feature & Description (Generated by GPT-4) & Interpretability Score \\ 
 \hline
 1-0000 & parts of individual names, especially last names. & 0.33 \\ 
 \hline
 1-0001 & actions performed by a subject or object. & -0.11 \\
 \hline
 1-0002 & instances of the letter `W' and words beginning with `w'. & 0.55 \\ 
 \hline
 1-0003 & the number `5' and also records moderate to low activation for personal names and some nouns. & 0.57 \\
 \hline
 1-0004 & legal terms and court case references. & 0.19 \\
 \hline
\end{tabular}
    \caption{Results of autointerpretation on the first five features found in the layer 1 residual stream. Autointerpretation produces a description of what the feature means and a score for how well that description predicts other activations.}
    \label{table:autointerpretation}
\end{table}

